\documentclass[11pt]{article}

\usepackage[T1]{fontenc}
\usepackage{amsmath,amssymb,amsthm}
\usepackage[margin=1in]{geometry}
\usepackage{xcolor}
\usepackage{hyperref}
\hypersetup{colorlinks=true,linkcolor=blue,urlcolor=blue,citecolor=blue}

\theoremstyle{plain}
\newtheorem{theorem}{Theorem}[section]
\newtheorem{lemma}[theorem]{Lemma}
\newtheorem{proposition}[theorem]{Proposition}
\newtheorem{corollary}[theorem]{Corollary}
\newtheorem{conjecture}[theorem]{Conjecture}
\theoremstyle{definition}
\newtheorem{definition}[theorem]{Definition}
\theoremstyle{remark}
\newtheorem{remark}[theorem]{Remark}

\newcommand{\R}{\mathbb R}
\newcommand{\C}{\mathbb C}
\newcommand{\Z}{\mathbb Z}
\newcommand{\Q}{\mathbb Q}
\newcommand{\CAZ}{\mathrm{CAZ}}
\newcommand{\dist}{\operatorname{dist}}
\newcommand{\re}{\operatorname{Re}}
\newcommand{\iC}{\iota_{C_\tau}}
\newcommand{\El}{\mathcal E_n}

\title{Convergence of iterated SOCP refinement for CAZAC feasibility}
\author{Jeffery Kline}
\date{Version 0.1.0, August 2026\\
\href{https://doi.org/10.5281/zenodo.21766714}{doi:10.5281/zenodo.21766714}}

\begin{document}
\maketitle

\begin{abstract}
Constant-amplitude zero-autocorrelation (CAZAC) sequences have constant
magnitude in time and flat Fourier magnitude.  We study an iterative
construction that relaxes these two nonconvex equality constraints to convex
inequalities and repeatedly solves a second-order cone program, using the
previous optimizer as the next objective direction.  At relaxation parameter
$\tau=1$, the CAZAC systems are exactly the maximal-norm points of the
resulting compact convex set.

For each fixed $\tau\ge1$, every nonzero trajectory converges to a single
fixed point and has finite total path length.  This conclusion follows by
applying a known
Kurdyka--{\L}ojasiewicz/Frank--Wolfe descent template to the CAZAC
norm-maximization formulation.  Global convergence is not global success:
we construct exposed non-CAZAC fixed points, so no deterministic choice of
optimizer can guarantee convergence to the CAZAC set.  Locally, a regularity
hypothesis gives a sharper conclusion.  Near a regular CAZAC limit the
constraint slack grows linearly, the relevant KL exponent is zero, and the
iteration reaches the limit exactly after finitely many steps.

We compute the regularity condition at every Zadoff--Chu tuple.  If
$n=\prod_p p^{a_p}$ and
$d_{\max}=\prod_p p^{\lfloor a_p/2\rfloor}$, then the active-gradient
dependency space has dimension $d_{\max}$, for odd and even $n$ alike.
Thus the local finite-arrival theorem applies at Zadoff--Chu limits exactly
when $n$ is squarefree.  Earlier constructions already supplied continuous
CAZAC families at non-squarefree lengths; the result here is the exact local
count and its connection to convergence.  We also apply the general
convergence argument to the elliptope and prove a partial symmetry-restricted
statement for Bj\"orck sequences; full Bj\"orck regularity is prior work.
The paper separates unconditional proofs, the regularity-conditional result,
numerical evidence, and open basin-of-attraction questions.
\end{abstract}

\section{Introduction}

A constant-amplitude zero-autocorrelation (CAZAC) sequence is a finite complex
sequence whose entries all have the same magnitude and whose discrete Fourier
transform also has constant magnitude.  Such sequences are useful in waveform
design, synchronization, and coding, and they form a nonconvex
feasibility set: one must satisfy a torus constraint in time and another in
frequency at the same time.  Zadoff--Chu chirps give explicit examples at
every length \cite{Chu72}, but the broader geometry of the feasible set and
the behavior of iterative construction methods remain difficult.

This paper studies a convex-relaxation method for that feasibility problem.
For $N$ coupled length-$n$ sequences, let $C_\tau$ impose upper bounds on
each time-domain amplitude and on the coupled Fourier energy at each
frequency.  The set $C_\tau$ is compact, convex, semialgebraic, and
second-order-cone representable.  At $\tau=1$, its points of largest squared
norm are exactly the coupled CAZAC systems.  Starting from a nonzero vector,
we repeatedly maximize the linear functional defined by the current iterate:
\[
 x^{\ell+1}\in\operatorname*{arg\,max}_{x\in C_\tau}
       \langle x^\ell,x\rangle_\R.
\]
This is iterated linear optimization, or a generalized power method, and each
step is an SOCP.

The first result separates convergence of the iteration from success of the
construction.  Theorem~\ref{thm:A} proves that every measurable optimizer
selection has finite total path length and converges to one fixed point.  The
descent mechanism is not new: it is the quadratic-objective instance of the
Kurdyka--{\L}ojasiewicz/Frank--Wolfe framework developed in greater generality
by Akta\c{s} and Kroer \cite{AK25}, building on standard KL descent theory
\cite{ABS13}.  The CAZAC-specific contribution is the norm-maximization
reduction that makes this framework apply.  Theorem~\ref{thm:B} then shows why
convergence alone cannot guarantee success: exposed non-CAZAC fixed
points exist, so a deterministic selection rule cannot guarantee arrival at
the target set.

The local theory is stronger.  Under a regularity condition on the active
constraint gradients, Theorem~\ref{thm:C} proves linear growth of the
constraint slack, KL exponent zero, and exact arrival at the limiting CAZAC
point after finitely many steps.  Theorem~\ref{thm:D} determines when this
hypothesis holds at Zadoff--Chu tuples.  If
$n=\prod_p p^{a_p}$, the dependency space has dimension
$d_{\max}=\prod_p p^{\lfloor a_p/2\rfloor}$; hence regularity holds exactly
for squarefree $n$.  The proof uses closure of quadratic chirps under the
DFT to replace a rank computation by the congruence $t^2\equiv0\pmod n$.

Earlier CAZAC work already shows why degeneracy can occur.  Backelin
\cite{Backelin89}, Popovi\'c \cite{Popovic92}, and Bj\"orck--Saffari
\cite{BjorckSaffari95} already constructed continuous families at
non-squarefree lengths and therefore anticipated the existence of the
degeneracy.  Theorem~\ref{thm:D} supplies the exact count at every
Zadoff--Chu tuple, for both parities, and connects it to the convergence
rate.  The paper does not identify its dependency space with those earlier
families, and it does not claim that the existence of non-squarefree
degeneracy is new.

The same method also applies to the elliptope.  The global
finite-length argument transfers to the elliptope of correlation matrices,
where it proves single-point convergence even though the fixed-point set is
positive-dimensional for $n>3$.  The pointwise attractive-fixed-point
dichotomy follows from Felzenszwalb, Klivans, and Paul \cite{FKP21}, but it
remains open whether almost every initialization reaches a rank-one point.
For Bj\"orck's Legendre-symbol construction, we prove regularity only
on a symmetry-invariant subspace; full regularity is proved by Benoist
\cite{Benoist24} and is not a contribution of this paper.

For $N=1$ and odd prime $n$, Benoist \cite[Prop.~A.2]{Benoist24} proves
projective transversality at Gaussian chirps, the prime-length counterpart
of hypothesis (R).  Theorem~\ref{thm:D} supplies the exact dependency count
and extends the calculation to every length, both parities, every coprime
root, and coupled systems with $N\ge1$.

The next section defines the relaxation and its fixed points.  Theorem A and
the elliptope application follow, then the obstruction theorem, the
regularity-conditional local result, and the Zadoff--Chu dependency
calculation.  The paper closes with open problems and a claim-by-claim
verification ledger.  Numerical experiments support selected calculations
and basin questions, but they are kept distinct from the proofs.

\section{Setup}\label{sec:setup}

Fix $n\ge1$ and $N\ge1$. Let $F$ denote the unnormalized discrete Fourier
transform on $\C^n$, $F(u)(k)=\sum_m u(m)e^{-2\pi ikm/n}$, so Parseval's
identity reads $\|Fu\|^2=n\|u\|^2$. The state space is
$x=(x_0,\dots,x_{N-1})\in(\C^n)^N\cong\R^{2Nn}$, with real inner product
$\langle u,v\rangle_\R=\re\langle u,v\rangle$ and induced norm
$\|x\|^2=\sum_{j,k}|x_j(k)|^2$. Write the \emph{coupled spectrum}
$S_x(k)=\sum_{j=0}^{N-1}|F(x_j)(k)|^2$.

\begin{definition}[Coupled CAZAC system]
$x$ is a coupled CAZAC system if $|x_j(k)|=1$ for all $j,k$ and
$S_x(k)=Nn$ for all $k$. Write $\CAZ(n,N)$ for this set; for $N=1$ this is
the classical CAZAC condition. $\CAZ(n,N)$ is nonempty for every $n\ge1$,
$N\ge1$: a Zadoff--Chu sequence exists for every length $n$ \cite{Chu72},
and $N$ copies of one satisfy the coupled constraint.
\end{definition}

For $\tau\ge1$, the convex relaxation is
$$C_\tau=\bigl\{x\in(\C^n)^N: |x_j(k)|\le\tau\ \forall j,k,\ \ S_x(k)\le Nn\ \forall k\bigr\}.$$
$C_\tau$ is compact (amplitude bounds), convex, semialgebraic, and
second-order-cone representable (each constraint is a Euclidean-norm bound
on a linear image of $x$); it is invariant under global phase rotation of
each $x_j$, cyclic shifts, and (for the coupled constraint) any
transformation preserving $S_x$.

\textbf{Iteration.} Given $x^0\ne0$, the algorithm is \emph{iterated linear
optimization}:
\begin{equation}
x^{\ell+1}\in\operatorname*{arg\,max}_{x\in C_\tau}\langle x^\ell,x\rangle_\R. \tag{$\star$}
\end{equation}
The argmax may be non-unique; $(\star)$ permits any measurable selection.
The reference implementation performs ten preliminary solves with
$\tau=1/\ell$ for $\ell=1,\dots,10$, then resets to $\tau=1$ for the main
phase.  All theorems below concern an iteration at one fixed $\tau\ge1$
(in particular $\tau=1$), and hold for any measurable selection rule.  The
preliminary solves merely choose the input direction for the analyzed main
phase; they are not part of any theorem.

\begin{definition}[Fixed point]
$x^*$ is a fixed point of $(\star)$ if
$x^*\in\arg\max_{x\in C_\tau}\langle x^*,x\rangle_\R$, equivalently
$\sigma_{C_\tau}(x^*)=\|x^*\|^2$, where
$\sigma_C(r)=\max_{x\in C}\langle r,x\rangle_\R$ is the support function of $C$.
\end{definition}

Three increasingly strong convergence questions organize everything that
follows: (1) does $x^\ell$ converge to a single point, not merely have
limit points; (2) is that limit CAZAC; (3) at what rate.

\section{The iteration as a generalized power method}

\begin{proposition}\label{prop:1}
Let $f(x)=\tfrac12\|x\|^2$, so $f$ is convex and $\nabla f(x)=x$. Then
$(\star)$ is exactly
\[
x^{\ell+1}\in\arg\max_{x\in C_\tau}
    \langle\nabla f(x^\ell),x\rangle_\R.
\]
This is the conditional-gradient or generalized-power-method update for
maximizing $f$ over the compact convex set $C_\tau$ \cite{JNRS10}, also
called iterated linear optimization \cite{FKP21}.
\end{proposition}
\begin{proof}
Immediate from the definitions.
\end{proof}

\begin{proposition}[Max-norm points are exactly CAZAC points]\label{prop:2}
For every $x\in C_\tau$ ($\tau\ge1$), $\|x\|^2\le Nn$. For $\tau=1$,
equality holds iff $x\in\CAZ(n,N)$. The maximum is attained for every
$n,N,\tau$.
\end{proposition}
\begin{proof}
By Parseval, $\|x\|^2=\sum_j\|x_j\|^2=\tfrac1n\sum_k S_x(k)\le\tfrac1n\cdot n\cdot Nn=Nn$,
with equality iff $S_x(k)=Nn$ for every $k$. Separately, at $\tau=1$,
$\|x\|^2=\sum_{j,k}|x_j(k)|^2\le Nn$ with equality iff every $|x_j(k)|=1$.
Both chains tight simultaneously forces both CAZAC conditions; conversely a
CAZAC system attains $Nn$. Attainment: $N$ copies of a Zadoff--Chu sequence
of length $n$ lie in $C_1\subseteq C_\tau$ with norm$^2=Nn$.
\end{proof}

This reframes the nonconvex feasibility problem ``find a CAZAC system'' as
the convex maximization $\max_{x\in C_1}\|x\|^2$, whose optimal value $Nn$
is known a priori: an iterate is $\delta$-suboptimal iff $Nn-\|x^\ell\|^2=\delta$.

\begin{proposition}[Monotone ascent, square-summable increments]\label{prop:3}
Along the iteration at fixed $\tau$, with $x^\ell\ne0$ and \textbf{$x^0\in
C_\tau$}:
\begin{enumerate}
\item[\rm(1)] $\langle x^\ell,x^{\ell+1}\rangle_\R\ge\|x^\ell\|^2$;
\item[\rm(2)] $\|x^{\ell+1}\|\ge\|x^\ell\|$, and $\|x^\ell\|\uparrow\rho\le\sqrt{Nn}$;
\item[\rm(3)] $\|x^{\ell+1}-x^\ell\|^2\le\|x^{\ell+1}\|^2-\|x^\ell\|^2$, hence
$\sum_{\ell\ge0}\|x^{\ell+1}-x^\ell\|^2\le Nn-\|x^0\|^2$ and
$\|x^{\ell+1}-x^\ell\|\to0$.
\end{enumerate}
\end{proposition}
\begin{proof}
Parts (1)--(2) and the bound in (3) instantiate the argument of
\cite{FKP21}'s Theorem~1 for this $C_\tau$; only the explicit telescoped
constant $Nn-\|x^0\|^2$ in (3) is not in their statement, since it uses
Prop.~\ref{prop:2}'s CAZAC-specific bound $Nn$ in place of a general
compactness bound.
(1) $x^\ell\in C_\tau$ is feasible for the SOCP whose optimizer is
$x^{\ell+1}$, so the optimum is at least $\langle x^\ell,x^\ell\rangle_\R=\|x^\ell\|^2$;
this is where $x^\ell\in C_\tau$ is used, which is why the hypothesis is
stated for $x^0$ and not just $x^{\ell\ge1}$ (every $x^{\ell+1}$ for
$\ell\ge0$ is automatically in $C_\tau$, being an argmax over $C_\tau$).
(2) Cauchy--Schwarz: $\|x^\ell\|\|x^{\ell+1}\|\ge\langle x^\ell,x^{\ell+1}\rangle_\R\ge\|x^\ell\|^2$;
divide by $\|x^\ell\|$; monotone and bounded above by Prop.~\ref{prop:2}.
(3) $\|x^{\ell+1}-x^\ell\|^2=\|x^{\ell+1}\|^2+\|x^\ell\|^2-2\langle x^\ell,x^{\ell+1}\rangle_\R\le\|x^{\ell+1}\|^2-\|x^\ell\|^2$
by (1); telescope and apply (2).
\end{proof}

The hypothesis $x^0\in C_\tau$ costs nothing.  Whatever input direction is
handed to a fixed-$\tau$ phase, its first optimizer lies in $C_\tau$ by
construction.  Re-indexing that phase from its first optimizer therefore
restores the hypothesis with no loss; every result below that depends on
Prop.~\ref{prop:3} inherits this same one-step re-indexing.

Square-summable increments alone do not imply convergence of the full
sequence (harmonic-type sums can wander); closing exactly this gap is
Theorem A.

\begin{proposition}[Limit points are fixed points; CAZAC points are fixed points]\label{prop:4}
Every limit point of $(x^\ell)$ is a fixed point of $(\star)$. Every
$x^*\in\CAZ(n,N)$ is a fixed point of the $\tau=1$ iteration, and moreover
a global maximizer of its own functional:
$\arg\max_{x\in C_1}\langle x^*,x\rangle_\R\ni x^*$ with optimal value $Nn$.
\end{proposition}
\begin{proof}
Let $x^{\ell_j}\to\bar x$. By Prop.~\ref{prop:3}(3), $x^{\ell_j+1}\to\bar x$
as well. $\sigma_{C_\tau}$ is continuous (finite convex on $\R^{2Nn}$), and
$\langle x^{\ell_j},x^{\ell_j+1}\rangle_\R=\sigma_{C_\tau}(x^{\ell_j})$ by
definition of the update; pass to the limit. For the second claim, any
$y\in C_1$ satisfies $\langle x^*,y\rangle_\R\le\|x^*\|\|y\|\le\sqrt{Nn}\sqrt{Nn}=Nn=\|x^*\|^2$
by Cauchy--Schwarz and Prop.~\ref{prop:2}, with equality at $y=x^*$.
\end{proof}

The same Cauchy--Schwarz argument shows any maximal-norm point of \emph{any}
compact convex set is a fixed point of iterated linear optimization;
nothing CAZAC-specific is used beyond Prop.~\ref{prop:2} itself --- a fact
that drives the elliptope corollary in \S\ref{sec:elliptope}.

\section{Theorem A --- global convergence}

\begin{theorem}[Theorem A --- Proved, unconditional]\label{thm:A}
Fix $\tau\ge1$. For any $x^0\ne0$ and any measurable selection
$x^{\ell+1}\in\arg\max_{x\in C_\tau}\langle x^\ell,x\rangle_\R$, the sequence
$(x^\ell)$ has finite length, $\sum_\ell\|x^{\ell+1}-x^\ell\|<\infty$, and
converges to a single fixed point $\bar x\in C_\tau$.
\end{theorem}
\begin{proof}
Apply the abstract descent theorem of Attouch, Bolte, and Svaiter
\cite[Thm.~2.9]{ABS13} to $F(x)=-\tfrac12\|x\|^2+\iC$ (so minimizing $F$ is
maximizing $\|x\|^2/2$ on $C_\tau$):
\begin{itemize}
\item \textbf{Sufficient decrease} ($a=\tfrac12$): exactly Prop.~\ref{prop:3}(3).
\item \textbf{Relative error} ($b=1$): first-order
optimality of the argmax step gives $x^\ell\in N_{C_\tau}(x^{\ell+1})$, the
normal cone at the new iterate. Since $\partial F(x^{\ell+1})=-x^{\ell+1}+N_{C_\tau}(x^{\ell+1})$
(sum rule for $C^1$ plus indicator of a convex set), the vector
$w=x^\ell-x^{\ell+1}$ lies in $\partial F(x^{\ell+1})$, and
$\|w\|=\|x^{\ell+1}-x^\ell\|$ \emph{exactly}: the subgradient residual is
literally the step. This identity is specific to $f=\tfrac12\|x\|^2$; for a
generic convex objective the relative-error condition would be a genuine
additional hypothesis, but here it is free. (The identical construction,
for a general smooth strongly convex objective where it holds only as a
Lipschitz-constant inequality rather than an exact equality, is
Aktaş--Kroer's \cite{AK25} Lemma 3.1 eq.~(27)--(28); our case is the
special instance where that Lipschitz constant is exactly $1$, so the
inequality becomes an equality --- this observation is a corollary of
their general lemma, not an independent discovery.)
\item \textbf{Compactness/continuity}: iterates lie in compact $C_\tau$;
$F$ is continuous there.
\item \textbf{KL property}: $C_\tau$ is defined by polynomial (in)equalities
and $F$ is polynomial plus the indicator of a semialgebraic set, hence
semialgebraic, hence KL \cite{BDLS07}.
\end{itemize}
Critical points of $F$ are exactly the fixed points of $(\star)$: $0\in\partial F(x^*)$
iff $x^*\in N_{C_\tau}(x^*)$ (the sum rule above) iff $x^*$ maximizes
$\langle x^*,\cdot\rangle_\R$ over $C_\tau$, which is the Fixed-point
Definition of Section~\ref{sec:setup}. The theorem then gives finite length and
convergence to a single critical point.
\end{proof}

The theorem is silent on \emph{which} fixed point the limit is; that is
exactly where Theorem B enters. Selection-independence, and the fact that
the presolve is finitely many extra solves, mean the conclusion applies
verbatim to any concrete implementation of $(\star)$.

\begin{remark}[Novelty check against existing CAZAC-construction
algorithms]
The closest prior iterative CAZAC-construction method is IPUC (Iterative
Projection onto Unit Circle) \cite{Amis25}: alternate projecting onto the
unit-modulus constraint in time and in frequency, from a random seed.
IPUC's own authors report, at length $n=50$, $8$ of $20$ trajectories
with discrepancy frozen above $10^{-3}$ after $10^5$ iterations, and
state plainly that ``the analysis of this iterative function is beyond
the scope of this paper''; their practical remedy is to restart from a
new random seed rather than continue iterating. Theorem~A gives what
IPUC's own presentation explicitly does not attempt: an unconditional
proof that (a different, convex-relaxation-based) iteration converges,
from every initialization, to a single fixed point, with finite total
path length. The two algorithms are not the same iteration --- IPUC
projects directly onto the nonconvex unimodular constraint in both
domains, while $(\star)$ relaxes to the convex body $C_\tau$ --- so this
is not a convergence proof for IPUC itself, but it is, as far as we have
found, the first provable convergence guarantee for an iterative
CAZAC-construction method in this literature.
\end{remark}

\subsection{A second worked example: the elliptope}\label{sec:elliptope}

The proof of Theorem~\ref{thm:A} uses nothing about $C_\tau$ beyond
compactness and convexity, except at the KL step, which needs
semialgebraicity. Stated in general: for \emph{any} compact convex
$C\subset\R^d$ and the iteration
$x^{\ell+1}\in\arg\max_{x\in C}\langle x^\ell,x\rangle$, sufficient decrease
(Prop.~\ref{prop:3}(3)'s proof uses only feasibility of $x^\ell$ for the
next problem and Cauchy--Schwarz) and the relative-error identity (the
normal-cone argument above, using only $x^\ell\in N_C(x^{\ell+1})$ and
$\nabla(\tfrac12\|x\|^2)=x$) hold unconditionally; compactness gives
continuity. If $C$ is also semialgebraic, KL holds \cite{BDLS07} and
Theorem A's conclusion --- convergence to a single fixed point, not merely a
connected limit set --- follows for that $C$ with no CAZAC-specific input.

We exercise this on a genuinely different $C$: the \textbf{elliptope}
$$\El=\{X\in\R^{n\times n}: X=X^\top,\ X\succeq0,\ X_{ii}=1\ \forall i\},$$
the set of $n\times n$ correlation matrices, with iteration
$X^{\ell+1}\in\arg\max_{X\in\El}\langle X^\ell,X\rangle$
($\langle A,B\rangle=\operatorname{tr}(AB)$). This is Felzenszwalb, Klivans,
and Paul's own example \cite{FKP21}: a case where they concede the
fixed-point set is a continuum for $n>3$, and where their own worked
instance stops at set-level convergence.

\begin{proposition}[Proved]\label{prop:ell1}
$\El$ is compact, convex, and semialgebraic.
\end{proposition}
\begin{proof}
Convexity of the constraints is immediate. Compactness: $X\succeq0$ with
$X_{ii}=1$ gives $|X_{ij}|\le\sqrt{X_{ii}X_{jj}}=1$ by Cauchy--Schwarz on
the associated Gram vectors, so $\El$ is bounded, and closed as an
intersection of closed sets, hence compact. Semialgebraicity: for
symmetric $X$, $X\succeq0$ is equivalent by Sylvester's criterion to the
finite conjunction of polynomial inequalities ``every principal minor of
$X$ is $\ge0$''; the diagonal constraints are polynomial equalities. So
$\El$ is a basic semialgebraic set (a spectrahedron).
\end{proof}

\begin{proposition}[Proved]\label{prop:ell2}
Sufficient decrease, relative error, and compactness/continuity transfer
verbatim from Prop.~\ref{prop:3}(3) and the normal-cone argument of
Theorem~\ref{thm:A}, replacing ``$C_\tau$'' by ``$\El$'' and
$\langle\cdot,\cdot\rangle_\R$ by $\operatorname{tr}(\cdot\,\cdot)$.
\end{proposition}
\begin{proof}
Neither source proof uses any CAZAC structure; both use only convexity,
Cauchy--Schwarz, and the normal-cone characterization of the argmax, all of
which hold verbatim on $\El\subset\R^{\binom{n+1}2}$.
\end{proof}

\begin{proposition}[Proved, $n>3$]\label{prop:ell3}
For $n>3$, let
\[
\mathcal F=\{X\in\El:X\in\arg\max_{Y\in\El}\langle X,Y\rangle\}
\]
be the fixed-point set of the elliptope iteration. It is infinite, hence (being
semialgebraic) a genuine positive-dimensional continuum, not a singleton
or a finite set. The restriction $n>3$ is necessary: at
$n=2$, $\El=\{[\begin{smallmatrix}1&r\\r&1\end{smallmatrix}]:|r|\le1\}$ and
one checks directly that $\mathcal F=\{r\in\{-1,0,1\}\}$ is finite, so the
statement is false without $n>3$.
\end{proposition}
\begin{proof}
$\mathcal F$ is semialgebraic: $\El$ is semialgebraic
(Prop.~\ref{prop:ell1}) and $X\mapsto\sigma_{\El}(X):=\max_{Y\in\El}\langle
X,Y\rangle$ is semialgebraic by the Tarski--Seidenberg theorem (a
projection of a semialgebraic set), so $\mathcal F=\{X\in\El:\langle
X,X\rangle=\sigma_{\El}(X)\}$ is cut out by polynomial (in)equalities. A
semialgebraic subset of $\R^{\binom{n+1}2}$ that is infinite cannot be
$0$-dimensional (an o-minimal set with finitely many points in every cell
is finite), so it suffices to exhibit infinitely many fixed points.
Felzenszwalb, Klivans, and Paul \cite{FKP21} state exactly this: $\mathcal
L_n$ (their name for $\El$) has infinitely many fixed points for $n>3$ ---
this, not the example below, is where the hypothesis $n>3$ enters and
cannot be dropped. Concretely, $X^*=I_n$ is one such fixed point (for any
$n$), and its optimal face is
degenerate in the strongest possible sense: $\operatorname{tr}(I_nX)=\sum_iX_{ii}=n$
identically on $\El$, so the linear functional $\langle I_n,\cdot\rangle$
is constant there and $\arg\max_{X\in\El}\langle I_n,X\rangle=\El$ itself
--- every feasible point, including every rank-$1$ correlation matrix
$vv^\top$, is simultaneously optimal for $I_n$'s own direction, which in
particular makes $I_n$ a fixed point with a machine-exact identity. This
illustrates the phenomenon but is not itself the infinitude argument:
membership of $vv^\top$ or other points in $\arg\max\langle I_n,\cdot\rangle$
does not make them fixed points of the map, since a fixed point must be
optimal for \emph{its own} direction, not $I_n$'s; the infinitude of
$\mathcal F$ rests on \cite{FKP21}'s cited count, not on this example
alone.
\end{proof}

\begin{proposition}[Verified numerically, $n=5$]\label{prop:ell4}
Running the iteration on $\mathcal E_5$ from four starting points ---
$X^0=I_5$; $I_5$ plus a small random perturbation renormalized to unit
diagonal; and two independent random correlation matrices --- using an
interior-point SDP solver (CLARABEL via cvxpy), all four trajectories show
$\|X^{\ell+1}-X^\ell\|_F\to0$ (to $\sim10^{-15}$ by iteration $15$) and
settle on a single matrix. Three of the four escape the fixed-point
continuum entirely and converge to a rank-1 correlation matrix $vv^\top$
(eigenvalues $(0,0,0,0,5)$ to machine precision); the trajectory started at
$X^0=I_5$ is an edge case in which the solver's own deterministic tie-break
returns $I_5$ again, a valid fixed point, not a counterexample. Script:
\path{experiments/verify_elliptope_convergence.py} in the companion
repository \cite{CazacAlgorithmRepo}. A broader sweep (same repository,
\path{experiments/verify_elliptope_rank_sweep.py}) of $140$ random and
Gaussian-perturbed trials across $n\in\{3,4,5,8,10,15,20\}$ finds
\textbf{every trajectory converges to rank $1$}, no exceptions. This
remains \emph{verified numerically only}, at these tested $n$ and trial
counts, not a proof for every $n$ or every trajectory.
\end{proposition}

\begin{theorem}[Elliptope convergence --- tiers as stated]\label{thm:ell}
Because $\El$ is compact, convex, and semialgebraic (Prop.~\ref{prop:ell1}),
Theorem A's general form applies to it directly, for every $n$: the
iterated-linear-optimization sequence on the elliptope converges to a
single fixed point for every initialization --- \textbf{Proved}. For $n>3$
this holds \textbf{despite the fixed-point set being an infinite
continuum} (Prop.~\ref{prop:ell3}; for $n\le3$ the fixed-point set is
finite, so the continuum caveat is vacuous there but convergence is still
unconditional, by the same general argument). Which
fixed point is reached is settled only partially: the vertices (rank-$1$
points) are exactly the fixed points whose full neighborhood converges to
them, and no other fixed point has this property, \textbf{by direct
citation} to \cite{FKP21}'s own Theorem~20 (Remark~\ref{rem:fkp-thm20}
below). That generic nearby trajectories escape the continuum and settle
specifically onto a rank-$1$ point is \textbf{verified numerically} (140
trials, Prop.~\ref{prop:ell4}), not established for every $n$ or every
trajectory. Whether the measure of initializations converging elsewhere
is exactly zero is \textbf{open} (Conjecture, filed as Q7 below).
\end{theorem}

\begin{remark}[Relation to Felzenszwalb--Klivans--Paul]
FKP's Theorem~1 (limit points exist, are fixed points, form a connected
set) is the same starting point as Prop.~\ref{prop:4} here. From an
\emph{arbitrary} initialization, on a fixed-point set that is a
positive-dimensional continuum for $n>3$ (Prop.~\ref{prop:ell3}), their
only route to single-point convergence is an unnumbered remark that a
connected \emph{finite} fixed-point set is a singleton, which is
structurally
inapplicable here; a full read of the proof of their Theorem~1 confirms it
only ever derives connectedness and never separately invokes finiteness.
(FKP's own Theorem~20 does give single-point convergence from a
\emph{restricted} class of initializations --- any point within distance
$1$ of a vertex converges to it in one step, per
Remark~\ref{rem:fkp-thm20} below --- but this does not cover an arbitrary
starting point, which is the setting Theorem~\ref{thm:ell} needs.)
Theorem~\ref{thm:ell} closes exactly this gap via the KL/finite-length
route \cite{ABS13} rather than the connectedness-plus-finiteness route,
and is therefore not subsumed by \cite{FKP21} even in the
positive-dimensional-face case. This route itself, as noted at Theorem~A,
is not unique to this paper: Aktaş--Kroer's \cite{AK25} KL/Frank-Wolfe
framework (built on a different but interchangeable abstract descent
theorem, Bolte--Sabach--Teboulle--Vaisbourd) directly delivers the same
elliptope conclusion --- every hypothesis of their Theorem~3.2 (strongly
convex smooth objective, nonempty compact constraint set, KL) is verified
here by Prop.~\ref{prop:ell1}, so this is a second, independent route to
Theorem~\ref{thm:ell}, not merely a plausible one. \cite{AK25} does not
name FKP, does not run its framework on the elliptope itself under that
name (its own Max-Cut application factorizes to a product of spheres via
Burer--Monteiro, \S4.3), and does not discuss positive-dimensional
critical sets; the comparison to FKP should be read as ``a route FKP's own
paper does not take,'' not as ``a route no other paper could have taken.''
\end{remark}

\begin{remark}[FKP's own Theorem 20 bears directly on Q7]\label{rem:fkp-thm20}
Although \cite{FKP21} cannot prove Theorem~\ref{thm:ell}'s single-point
convergence claim, its own Theorem~20 answers a large part of Q7 (below)
by direct citation, without any new argument here: \emph{the vertices of
$\El$ (the rank-$1$ sign matrices $vv^\top$) are exactly the
\textbf{attractive} fixed points of the map $X\mapsto\arg\max_{Y\in\El}\langle
X,Y\rangle$}, in the pointwise sense that an entire neighborhood of a
vertex converges to it, while every non-vertex fixed point admits (their
Prop.~19) an explicit curve of nearby points whose norm strictly
increases, so it cannot attract a full neighborhood either. This is a
genuine dichotomy, not a numerical observation, and it is stronger than
what either paper's abstract convergence machinery gives alone: it says
which fixed points \emph{can} be limits of nearby trajectories, complementing
(not overlapping) Theorem~\ref{thm:ell}'s unconditional guarantee that
\emph{some} single limit exists. It falls short of Q7 as literally stated
for a precise reason: \cite{FKP21}'s Prop.~19 produces, in every
neighborhood of a non-vertex fixed point, \emph{at least one} nearby point
that escapes it --- exactly what ``not attractive'' (their Def.~6)
requires --- but this does not show \emph{every} nearby point escapes,
which is what ``repelling'' (their Def.~7) would additionally require.
Whether non-vertex elliptope fixed points are repelling outright, or can
be neither attractive nor repelling (as \cite{FKP21}'s own \S3, Example~5,
shows can happen in general iterated-linear-optimization systems, for an
unrelated three-dimensional cone, not the elliptope itself) is not settled
by Theorem~20 either way. If some non-vertex elliptope point is neither,
Theorem~20 alone would not rule out a positive-measure set of
initializations converging to it through a collective, not individually
attractive, mechanism. Q7's residual open content is exactly this
measure-zero refinement, structurally the same gap as Q6.
\end{remark}

\section{Theorem B --- obstruction: the limit need not be CAZAC}

\begin{theorem}[Theorem B --- Proved, unconditional]\phantomsection\label{thm:B}
\begin{enumerate}
\item[\rm(a)] The fixed points of $(\star)$ at $\tau=1$ with every
amplitude constraint slack are exactly the points whose \emph{coupled}
spectral profile satisfies $S_x(k)=Nn$ on some proper subset
$K\subsetneq\{0,\dots,n-1\}$ of frequencies and $S_x(k)=0$ off $K$ (i.e.\
$\widehat x_j(k)=0$ for every $j$, when $k\notin K$), with time-domain
modulus strictly less than $1$ everywhere; such a point has
$\|x\|^2=N|K|$. For $N=1$ this forces $|\widehat x(k)|=\sqrt n$ exactly on
$K$; for $N\ge2$ only the coupled sum is pinned, and the individual
$|\widehat x_j(k)|$ are otherwise free subject to
$\sum_j|\widehat x_j(k)|^2=Nn$ for $k\in K$.
\item[\rm(b)] There exist non-CAZAC fixed points that are the
\textbf{unique} maximizer of their own linear functional over $C_1$ ---
exposed points. For $n=2$: $x^*=(1,t)$, $t=\sqrt2-1$, is such a point; for
$N=2$, its diagonal doubling $((1,t),(1,t))$ is such a point.
\item[\rm(c)] Consequently no deterministic convergence-to-CAZAC theorem
is possible: every measurable selection rule in $(\star)$ retains the
point of (b) exactly, once reached.
\end{enumerate}
\end{theorem}
\begin{proof}[Proof of (a), general $N$]
By Prop.~\ref{prop:1}, $x^*\ne0$ is a fixed point iff it is a KKT point of
$\max_y\langle x^*,y\rangle_\R$ subject to $g_{j,m}(y)=|y_j(m)|^2-1\le0$
(each block $j=1,\dots,N$) and the single \emph{shared} constraint
$h_k(y)=S_y(k)-Nn\le0$; this is a convex program and Slater's condition
holds at $y=0$, so KKT is necessary and sufficient. The real gradients are
$\nabla_{y_j}g_{j,m}(y)=2y_j(m)e_m$ and, since $S_y(k)=\sum_j|\widehat
y_j(k)|^2$ couples all blocks, $\nabla_{y_j}h_k(y)=2F^*E_kFy_j$ for
\emph{every} $j$ with the \emph{same} multiplier $\mu_k$ --- this coupling
is the source of the difference from $N=1$. KKT at $x^*$ reads, per block,
$x^*_j=\Lambda_jx^*_j+F^*MFx^*_j$ for diagonal $\Lambda_j,M\ge0$
complementary to the active sets, with $M$ shared across $j$. If every
amplitude constraint is slack then each $\Lambda_j=0$, so
$x^*_j=F^*MFx^*_j$ for every $j$; applying $F$ and using $FF^*=nI$ gives
$\widehat x^*_j=nM\widehat x^*_j$ coordinatewise for every $j$, i.e.\ for
each $k$ either $\widehat x^*_j(k)=0$ for every $j$, or $\mu_k=1/n$. Let
$K=\{k:\exists j,\ \widehat x^*_j(k)\ne0\}\ne\varnothing$; for $k\in K$,
complementary slackness on the \emph{shared} constraint $h_k$ with
$\mu_k=1/n>0$ forces the active condition $S_x^*(k)=Nn$ --- the coupled
sum, not each block's own modulus, since $h_k$ is one constraint shared by
all $N$ blocks. Parseval (summed over blocks) gives
$\|x^*\|^2=\tfrac1n\sum_kS_{x^*}(k)=\tfrac1n|K|\cdot Nn=N|K|$; if $|K|=n$
this forces $\|x^*\|^2=Nn$, hence (Prop.~\ref{prop:2}) $x^*\in\CAZ(n,N)$,
contradicting amplitude-slackness. So $K\subsetneq\{0,\dots,n-1\}$. For
$N=1$, $S_x(k)=|\widehat x(k)|^2$, so the coupled condition $S_x(k)=n$
collapses to the per-block statement $|\widehat x(k)|=\sqrt n$; for
$N\ge2$ the individual $|\widehat x_j(k)|$ are unconstrained beyond
$\sum_j|\widehat x_j(k)|^2=Nn$ (e.g.\ $N=2,n=2$: $\widehat
x_0(0)=\sqrt3,\widehat x_1(0)=1$ satisfies $S_x(0)=4=Nn$ with neither
modulus equal to $\sqrt n=\sqrt2$).
Conversely, for such $x^*$ and any $y\in C_1$, Cauchy--Schwarz per
frequency gives
$\langle x^*,y\rangle_\R=\tfrac1n\sum_{k\in K}\sum_j\re(\widehat
x^*_j(k)\overline{\widehat y_j(k)})\le\tfrac1n\sum_{k\in
K}\sqrt{S_{x^*}(k)}\sqrt{S_y(k)}\le\tfrac1n\sum_{k\in K}\sqrt{Nn}\sqrt{Nn}=N|K|=\|x^*\|^2$,
attained at $y=x^*$ --- a direct support-function argument, no KKT needed
for sufficiency.
\end{proof}
\begin{proof}[Proof of (b), $n=2,N=1$]
For $y=(y_0,y_1)\in\C^2$ feasible in $C_1$ (so $|y_m|\le1$ for each $m$,
and $S_y(0)=|F(y)(0)|^2\le2$ directly, hence $|F(y)(0)|\le\sqrt2$),
$$\langle x^*,y\rangle_\R=(1-t)\,\re y(0)+t\,\re F(y)(0)\le(1-t)+t\sqrt2=1+t^2=\|x^*\|^2,$$
with equality iff $y(0)=1$ and $F(y)(0)=\sqrt2$, which forces
$y(1)=\sqrt2-1=t$, i.e.\ $y=x^*$ --- a genuine separating functional
exposing $x^*$ as a vertex-like point of $C_1$ even though
$x^*\notin\CAZ(2,1)$ ($\|x^*\|^2=4-2\sqrt2<2$). For the $N=2$ doubling
$((1,t),(1,t))$, the amplitude constraints are still per-block
($\re y_j(0)\le|y_j(0)|\le1$ for $j=0,1$), but the spectral constraint is
\emph{coupled} ($S_y(0)=\sum_j|F(y_j)(0)|^2\le4$, not two separate
per-block bounds of $2$), so the argument is not a coordinatewise repeat:
\begin{align*}
\langle x^*,y\rangle_\R
 &= (1-t)\bigl[\re y_0(0)+\re y_1(0)\bigr]
    +t\bigl[\re F(y_0)(0)+\re F(y_1)(0)\bigr]\\
 &\le 2(1-t)+t\sqrt2\sqrt{S_y(0)}\\
 &\le 2(1-t)+2\sqrt2\,t=2(1+t^2).
\end{align*}
Here the first bound uses the amplitude constraints directly, while
Cauchy--Schwarz gives
\[
\re F(y_0)(0)+\re F(y_1)(0)\le\sqrt2\sqrt{S_y(0)}
\]
on the coupled spectral term. The last expression matches
$\|x^*\|^2=2(1+t^2)$ for the
doubled point. Equality forces $\re y_j(0)=1$ for $j=0,1$ (each term of
the first bracket individually at its bound, so $y_j(0)=1$ exactly), and
in the Cauchy--Schwarz step: $\re F(y_0)(0)=\re F(y_1)(0)=|F(y_j)(0)|$
(parallel to $(1,1)$ and real, nonnegative) with $S_y(0)=4$ exactly,
giving $F(y_0)(0)=F(y_1)(0)=\sqrt2$; since $F(y_j)(0)=y_j(0)+y_j(1)=1+y_j(1)$,
this forces $y_j(1)=\sqrt2-1=t$ for both blocks, i.e.\ $y=x^*$ ---
again a genuine separating functional, uniquely exposing the doubled
point. Part (c) follows immediately from the existence of one such
retained point.
\end{proof}

\begin{remark}
Theorem A guarantees convergence to \emph{some} fixed point; Theorem B
shows the fixed-point set genuinely contains non-CAZAC points that trap
every selection rule. The remaining content of ``the algorithm works'' is a
probabilistic/basin-of-attraction statement (Q6 below), not a deterministic
one.
\end{remark}

\section{Theorem C --- finite exact termination at regular CAZAC limits}

Call $x^*\in\CAZ(n,N)$ \textbf{regular} (hypothesis (R)) if the only
linear dependency among the gradients of the active constraints at $x^*$
is the one forced by Parseval's identity; equivalently the dependency
space $K$ (Section~\ref{sec:zc} makes this precise) has $\dim K=1$. Let
$\sigma_s>0$ denote the associated smallest singular value at a regular
point.

\begin{lemma}[Quantitative Newton-type error bound]\label{lem:newton}
Let $\Phi:\R^d\to\R^q$ be a quadratic map (each component a polynomial of
degree $\le2$), $\Phi(x^*)=0$, $A:=D\Phi(x^*)$ surjective with smallest
singular value $\sigma>0$, and let $M$ be a Lipschitz constant for $D\Phi$
(finite and global, since $D\Phi$ is affine). Then for every $x$ with
$\|x-x^*\|\le\sigma/(8M)$ and $\|\Phi(x)\|\le\sigma^2/(16M)$,
$$\dist(x,\Phi^{-1}(0))\le\frac2\sigma\|\Phi(x)\|.$$
\end{lemma}
\begin{proof}
Since $\Phi$ is quadratic, Taylor's expansion is exact:
$\Phi(y+v)=\Phi(y)+D\Phi(y)v+B(v,v)$ with $B$ symmetric bilinear and
$\|B(v,v)\|=\tfrac12\|(D\Phi(y+v)-D\Phi(y))v\|\le\tfrac M2\|v\|^2$. Run the
modified Newton iteration $z_0=x$, $z_{i+1}=z_i-A^+\Phi(z_i)$, where $A^+$
is the Moore--Penrose pseudoinverse ($AA^+=I_q$ since $A$ is surjective,
$\|A^+\|=1/\sigma$). Write $b_i=\|\Phi(z_i)\|$, $d_i=-A^+\Phi(z_i)$, so
$\|d_i\|\le b_i/\sigma$. Then
$$\Phi(z_{i+1})=\Phi(z_i)-D\Phi(z_i)A^+\Phi(z_i)+B(d_i,d_i)=(A-D\Phi(z_i))A^+\Phi(z_i)+B(d_i,d_i),$$
so, using $\|A-D\Phi(z_i)\|\le M\|z_i-x^*\|$,
$$b_{i+1}\le\frac{M\|z_i-x^*\|}\sigma\,b_i+\frac M{2\sigma^2}b_i^2.\qquad(*)$$
\textbf{Claim:} every $z_i$ satisfies $\|z_i-x^*\|\le\sigma/(4M)$ and
$b_i\le b_02^{-i}$. By induction: given these, $(*)$ gives
$b_{i+1}\le\tfrac14b_i+\tfrac1{32}b_i<\tfrac12b_i$ (using
$b_i\le b_0\le\sigma^2/(16M)$). Summing the geometric increments,
$\sum_i\|d_i\|\le(b_0/\sigma)\sum_i2^{-i}=2b_0/\sigma\le\sigma/(8M)$, so
$\|z_{i+1}-x^*\|\le\|x-x^*\|+2b_0/\sigma\le\sigma/(4M)$, closing the
induction. The $z_i$ are Cauchy with limit $z_\infty$; $b_i\to0$ and
continuity of $\Phi$ give $\Phi(z_\infty)=0$, and $\|z_\infty-x\|\le2b_0/\sigma$.
No convexity is used anywhere in this argument.
\end{proof}

Fix $x^*\in\CAZ(n,N)$ and let $H=\{c\in\R^n:\sum_kc_k=0\}$, $P_H$ the
orthogonal projection onto $H$. Define the \textbf{balanced constraint map}
$\Phi_s(x)=\bigl(g(x),\,P_Hh(x)/n\bigr)\in\R^{Nn}\times H$, where
$g_{j,m}(x)=|x_j(m)|^2-1$ and $h_k(x)=S_x(k)-Nn$ (so $g(x)=0$ is the
amplitude constraint and $h(x)=0$ is the balanced spectral constraint).

\begin{lemma}[Zero set]\label{lem:phis-zero}
$\Phi_s^{-1}(0)=\CAZ(n,N)$, globally.
\end{lemma}
\begin{proof}
If $g(x)=0$ then $x$ is unimodular, so by Parseval
$\sum_kh_k(x)=n\|x\|^2-nNn=0$, i.e.\ $h(x)\in H$; then $P_Hh(x)=0$ forces
$h(x)=0$, so $x\in\CAZ(n,N)$. Conversely $\CAZ(n,N)\subseteq\{g=0,h=0\}\subseteq\Phi_s^{-1}(0)$.
\end{proof}

\begin{lemma}[Uniform derivative bounds]\label{lem:phis-lip}
The derivative map $D\Phi_s$ is globally $2\sqrt2$-Lipschitz. At any
$x^*\in\CAZ(n,N)$ it also satisfies
$\|D\Phi_s(x^*)\|_{\mathrm{op}}\le2\sqrt{N+1}$. Both bounds are independent
of $n$.
\end{lemma}
\begin{proof}
\emph{(Lipschitz.)} The rows of $Dg$ are $2x_j(m)$, supported on disjoint
real $2$-planes, so $\|Dg(x)-Dg(y)\|_F=2\|x-y\|$. The $h_k$-row of
$Dh/n$ in block $j$ is (the real form of) $2\hat x_j(k)f_k/n$
($f_k(m)=e^{2\pi ikm/n}$), so by Parseval
$\|Dh(x)/n-Dh(y)/n\|_F^2=4\|x-y\|^2$; $P_H$ is a contraction. Stacking:
Lipschitz constant $\le\sqrt{4+4}=2\sqrt2$.
\emph{(Operator norm at $x^*$.)} $Dg(x^*)Dg(x^*)^\top=4I$ (disjoint unit
$2$-planes, $|x^*_j(m)|=1$); the Gram matrix of the unprojected scaled
spectral rows is $4N\delta_{kk'}$ (Parseval per block, $S_{x^*}\equiv Nn$).
The two blocks have operator norms $2$ and $2\sqrt N$; the stack has norm
$\le\sqrt{4+4N}=2\sqrt{N+1}$.
\end{proof}

\textbf{Hypothesis (R) restated exactly.} $D\Phi_s(x^*)$ is surjective onto
$\R^{Nn}\times H$ iff $\dim K=1$, where $K$ is the dependency space of
Section~\ref{sec:zc}: a pair $(a,c)$ with $c\in H$ annihilates the row
space of $D\Phi_s(x^*)$ iff $\sum_{j,m}a_{j,m}Dg_{j,m}+\sum_k(c_k/n)Dh_k=0$,
and by the definition of $K$ these solutions are parametrized by $c\in K$
(with $a$ determined by $c$). The Parseval solution $c=(1,\dots,1)\notin H$,
so the annihilator inside $\R^{Nn}\times H$ has dimension $\dim(K\cap H)=\dim K-1$
($K$ always contains the constant vector, which is not in $H$); hence
surjectivity $\iff\dim K=1$, matching hypothesis (R) as stated above.
Write $\sigma_s>0$ for the smallest singular value of $D\Phi_s(x^*)$ when
(R) holds, and $\rho:=\sigma_s/(16\sqrt2)$ (so $\sigma_s/(8M)=\rho$ with
$M=2\sqrt2$, Lemma~\ref{lem:phis-lip}).

\begin{theorem}[Theorem C --- Proved conditional on (R)]\label{thm:C}
Near a regular $x^*\in\CAZ(n,N)$:
\begin{enumerate}
\item[\rm(a)] \emph{Sharp (linear) growth:}
$D(x)\ge\dfrac{\sigma_s}{2\sqrt2}\dist(x,\CAZ(n,N))$ for $x\in C_1$ near
$x^*$, where $D(x)=Nn-\|x\|^2$.
\item[\rm(b)] \emph{KL exponent $0$:}
$\dist(0,\partial F(x))\ge c_0:=\sigma_s/(8\sqrt2)$ for all non-CAZAC $x$
in a neighborhood of $x^*$.
\item[\rm(c)] \emph{Finite exact arrival:} if the Theorem~A limit $\bar x$
is a regular CAZAC point, then $x^\ell=\bar x$ \textbf{exactly} for all
sufficiently large $\ell$. The number of ``large'' steps before exact
arrival is at most $(Nn-\|x^0\|^2)/c_0^2$.
\end{enumerate}
\end{theorem}
\begin{proof}
(a) The objective deficit is simultaneously the $\ell^1$ slack of both
constraint families:
$D(x)=\sum_{j,m}(1-|x_j(m)|^2)=\tfrac1n\sum_k(Nn-S_x(k))$, all terms
$\ge0$. In particular,
\begin{align*}
\|\Phi_s(x)\|
 &=\sqrt{\|g(x)\|^2+\|P_Hh(x)/n\|^2}\\
 &\le\sqrt{D(x)^2+D(x)^2}=\sqrt2D(x).
\end{align*}
\emph{Case 1: $\sqrt2D(x)\le\sigma_s^2/(16M)$, $M=2\sqrt2$.} Apply
Lemma~\ref{lem:newton} to $\Phi=\Phi_s$ (quadratic; zero set $\CAZ(n,N)$ by
Lemma~\ref{lem:phis-zero}; surjective at $x^*$ by (R)):
$\dist(x,\CAZ(n,N))\le(2/\sigma_s)\|\Phi_s(x)\|\le(2\sqrt2/\sigma_s)D(x)$,
i.e.\ $D(x)\ge(\sigma_s/(2\sqrt2))\dist(x,\CAZ(n,N))$.
\emph{Case 2: $\sqrt2D(x)>\sigma_s^2/(16M)$}, i.e.\ $D(x)>\sigma_s^2/64$.
Since $x$ is in the stated neighborhood, $\dist(x,\CAZ(n,N))\le\|x-x^*\|\le\rho=\sigma_s/(16\sqrt2)$,
so $(\sigma_s/(2\sqrt2))\dist(x,\CAZ(n,N))\le\sigma_s^2/64<D(x)$: the bound
holds directly. Either case gives (a).

(b) Let $d:=\dist(x,\CAZ(n,N))>0$ and pick a nearest $z\in\CAZ(n,N)$
(exists, $\CAZ(n,N)$ compact). Any $w\in\partial F(x)$ has the form
$w=v-x$ with $v\in N_{C_1}(x)$ (the sum rule, as in Theorem~A's proof).
Since $z\in C_1$, the normal-cone inequality gives $\langle v,z-x\rangle_\R\le0$,
hence, using $\|z\|^2=Nn$,
$$\langle w,z-x\rangle_\R\le-\langle x,z-x\rangle_\R=\|x\|^2-\langle x,z\rangle_\R
=-\tfrac12(Nn-\|x\|^2)+\tfrac12d^2=-(F(x)-F(x^*))+\tfrac{d^2}2,$$
so Cauchy--Schwarz on the left gives $\|w\|\,d\ge(F(x)-F(x^*))-d^2/2$.
By (a) (in the form $F(x)-F(x^*)\ge(\sigma_s/(4\sqrt2))d$, i.e.\ (a) with
$D(x)=2(F(x)-F(x^*))$), dividing by $d$:
$\|w\|\ge(F(x)-F(x^*))/d-d/2\ge\sigma_s/(4\sqrt2)-d/2$. Since
$d\le\|x-x^*\|\le\rho=\sigma_s/(16\sqrt2)$, this gives
$\|w\|\ge\sigma_s/(4\sqrt2)-\sigma_s/(32\sqrt2)=(7/32)\sigma_s/\sqrt2\ge\sigma_s/(8\sqrt2)=c_0$.
Taking the infimum over $w\in\partial F(x)$ proves (b).

(c) \emph{(Full chain.)} Let $\bar x$ be the Theorem~A limit, regular with
constants $\rho,\sigma_s,c_0$ as in (a)/(b). Since $x^\ell\to\bar x$ and
$\|x^{\ell+1}-x^\ell\|\to0$ (Theorem~A's square-summability), fix $L_0$ so
that for all $\ell\ge L_0$: $x^{\ell+1}$ lies in the (b)-neighborhood of
$\bar x$, and $\|x^{\ell+1}-x^\ell\|<c_0$. By the exact identity from
Theorem~A's proof, $w^{\ell+1}:=x^\ell-x^{\ell+1}\in\partial F(x^{\ell+1})$
with $\|w^{\ell+1}\|=\|x^{\ell+1}-x^\ell\|$ \emph{exactly}, at every $\ell$
-- no approximation. Suppose toward contradiction $x^{\ell+1}\notin\CAZ(n,N)$
for some $\ell\ge L_0$. Then $x^{\ell+1}$ is a non-CAZAC point inside the
(b)-neighborhood, so (b) gives $\dist(0,\partial F(x^{\ell+1}))\ge c_0$. But
$w^{\ell+1}\in\partial F(x^{\ell+1})$ and $\|w^{\ell+1}\|<c_0$ by choice of
$L_0$ --- a contradiction, since $w^{\ell+1}$ itself witnesses a subgradient
of norm $<c_0$. Hence $x^{\ell+1}\in\CAZ(n,N)$ \textbf{exactly}: the
identity is exact and (b)'s bound is a hard floor on every non-CAZAC point
in the neighborhood, so a subgradient below that floor cannot occur unless
the point already lies on $\CAZ(n,N)$. Write $z:=x^{\ell+1}\in\CAZ(n,N)$.
The next step maximizes $\langle z,y\rangle_\R$ over $y\in C_1$; by
Cauchy--Schwarz and Prop.~\ref{prop:2} ($\|z\|=\sqrt{Nn}$ is the global
max-norm), $\langle z,y\rangle_\R\le\|z\|\|y\|\le Nn$ with equality iff
$y=z$ --- the argmax is the singleton $\{z\}$, so $x^{\ell+2}=z$, and by
induction $x^m=z$ for all $m>\ell$; in particular $\bar x=z\in\CAZ(n,N)$.
The step-count bound is Theorem~A's
$\sum_\ell\|x^{\ell+1}-x^\ell\|^2\le Nn-\|x^0\|^2$ telescoped against
$c_0^2$ per large step (increment $\ge c_0$); since every step after
arrival has increment exactly $0$ (the sequence is constant at $z$), every
large step occurs strictly before arrival, so this global count already
\emph{is} the count of large steps before arrival, with no separate
argument needed to exclude post-arrival steps.
\end{proof}

\begin{remark}
This is a genuine proof by contradiction, not a heuristic squeeze: an
exact identity (equality, from Theorem~A) is jointly satisfiable with a
hard lower bound (inequality, from (b)) on the same quantity only at
$\dist=0$. Parts (a), (b), (c) share one status tag because they rest on
exactly the same two inputs --- hypothesis (R) at $\bar x$ (giving
$\sigma_s,\rho$) and the exact H2 identity from Theorem~A's proof --- and
(c) adds no assumption beyond the theorem's own hypothesis
``$\bar x\in\CAZ(n,N)$'' plus the unconditional singleton-argmax fact.
\end{remark}

\begin{remark}[The $N=1$ vs.\ $N=2$ gap is local-rate-invariant]
Theorem C's constants do not distinguish $N=1$ from $N=2$: first-order
sharpness at a regular point decides the local geometry identically. The
empirically observed difference in success probability between the
single-sequence and coupled formulations is therefore a global
(basin-of-attraction) phenomenon, not a local-rate one.
\end{remark}

\begin{remark}[Novelty check against generic identifiability theory]
Neither the Lewis--Luke--Malick theory of local linear convergence for
averaged/alternating projections \cite{LLM09} nor the Drusvyatskiy--Lewis
generic identifiability program \cite{DL12} delivers
Theorem~C. The former's core hypothesis, strongly regular (transversal)
intersection, is exactly what Section~\ref{sec:zc} shows \emph{fails} at
every CAZAC point (the universal Parseval dependency), and even where
satisfiable its conclusion is R-linear convergence of a different
algorithm, never finite exact arrival. The latter's genericity condition is
plausibly satisfied at CAZAC points, but its machinery yields only a
qualitative locally-minimal-identifiable-set statement, not a quantitative
sharp-growth constant, KL exponent, or finite-step termination guarantee.
\end{remark}

\section{Theorem D and regularity (R)}\label{sec:zc}

Write $x_u(m)=\omega^{-uT_m}$ ($\omega=e^{2\pi i/n}$, $T_m=m(m+1)/2$, $n$
odd) or $x_u(m)=\omega_{2n}^{-um^2}$ ($\omega_{2n}=e^{i\pi/n}$, $n$ even)
for the Zadoff--Chu sequence of root $u$, $\gcd(u,n)=1$ \cite{Chu72}. At a
CAZAC tuple all constraints are active; write
$$K=\{c\in\R^n: \overline{x_j(m)}\,\bigl(F^*(c\odot\widehat x_j)\bigr)(m)\in\R\ \ \forall j,m\}$$
for the active-gradient dependency space (the tangent space of the
constraint intersection has dimension $(N-1)n+\dim K$; hypothesis (R) is
$\dim K=1$).

\begin{theorem}[Theorem D --- Proved for all $n\ge1$, odd and even]\label{thm:D}
Let $n=\prod_p p^{a_p}$, $d_{\max}=\prod_p p^{\lfloor a_p/2\rfloor}$ (the
largest $d$ with $d^2\mid n$). Then for every $n\ge1$, every $N\ge1$, every
tuple of coprime roots $u_1,\dots,u_N$,
$$K=\{c\in\R^n: c(k+d_{\max})=c(k)\ \forall k\},\qquad \dim_\R K=d_{\max},$$
independent of $N$, of the roots, and of the parity of $n$.
\end{theorem}
\begin{proof}
\textbf{Step 0.} A dependency $\sum_{j,m}a_{j,m}\nabla g_{j,m}+\sum_kc_k\nabla h_k=0$
evaluated on a perturbation is equivalent to $F^*(c\odot\widehat x_j)$
being coordinatewise a real multiple of $x_j$ --- the defining condition
of $K$; dependencies with $c=0$ force $a=0$ since amplitude gradients have
disjoint supports. So the dependency space is (isomorphic to) $K$.

\textbf{Step 1 (self-transform).} For odd $n$, writing $k=us$,
$\widehat x_u(k)=G_u\,\omega^{uT_{u^{-1}k}}$ with $G_u=\sum_j\omega^{-uT_j}$,
$|G_u|^2=n$ by Parseval, using the exact identity $T_{m+s}=T_m+T_s+ms$. For
even $n$, the parallel identity holds with $T_m$ replaced by $m^2$,
exponents tracked mod $2n$: $\widehat x_u(k)=G_u\,\omega_{2n}^{us^2}$,
using $\omega_{2n}^{v(j+n)^2}=\omega_{2n}^{vj^2}$ (an exact integer
identity valid because $n$ is even).

\textbf{Step 2 (reduction to a correlation equation).} Fix a block with
root $u$ and write $q_m=\overline{x_u(m)}\,(F^*(c\odot\widehat x_u))(m)$.
Substituting Step 1, reindexing $k=us\bmod n$ (a bijection; $c'_s:=c_{us}$
again real), and using the same addition-law identity as Step 1 gives, in
both parities, $q_m=G_u\sum_sc'_s\,\phi(s+m)$, where $\phi(j)=\omega^{uT_j}$
(odd $n$) or $\phi(j)=\omega_{2n}^{uj^2}$ (even $n$). Since $c'$ is real,
$q_m\in\R$ for all $m$ iff $q_m-\overline{q_m}=0$:
$$\sum_sc'_s\,e(s+m)=0\quad\forall m,\qquad e(j):=G_u\,\phi(j)-\overline{G_u}\,\overline{\phi(j)}$$
--- the vanishing of a convolution against the \emph{reality-defect}
kernel $e$, not mere reality
of $\sum_sc'_s\phi(s+m)$ itself, which would be a strictly stronger, false
requirement: $\phi$ alone is not conjugate-symmetric, and $\hat\phi$ never
vanishes (it is a pure Gauss sum of constant modulus $\sqrt n$), so that
reading would force $\hat c'\equiv0$, contradicting $\dim K\ge1$.

\textbf{Step 3 (Fourier solve).} Taking the length-$n$ DFT in $m$ turns the
convolution into $\hat e(t)\,\widehat{c'}(-t)=0$ for every $t$, i.e.\
$\hat c'(t)=0$ wherever $\hat e(-t)\ne0$. Writing $\sigma=u^{-1}t$ and using
Step~1's formula for $\phi$ with $|G_u|^2=n$:
$\hat e(t)=n\bigl(\omega^{u(\sigma-T_\sigma)}-\omega^{uT_\sigma}\bigr)$
(odd $n$, using $T_{-\sigma}=T_\sigma-\sigma$) or
$\hat e(t)=n\bigl(\omega_{2n}^{-u\sigma^2}-\omega_{2n}^{u\sigma^2}\bigr)$
(even $n$). In both parities $\hat e(t)=0$ exactly when
$\sigma^2\equiv0\pmod n$ (odd: divide the exponent congruence by $u$ to get
$\sigma-T_\sigma\equiv T_\sigma$, i.e.\ $\sigma\equiv2T_\sigma=\sigma^2+\sigma$;
even: $2u\sigma^2\equiv0\pmod{2n}\iff u\sigma^2\equiv0\pmod n$, and $u$ is a
unit): the zero set is the subgroup $Z_0=\{t: t^2\equiv0\bmod n\}=m_0\Z_n$,
$|Z_0|=n/m_0=d_{\max}$ where $m_0=\prod_pp^{\lceil a_p/2\rceil}$.

\textbf{Step 4.} $\hat c'(t)=\hat c(u^{-1}t)$, so ``$\operatorname{supp}\hat c'\subseteq Z_0$''
is the same condition for every root and block, giving
$K=\{c:\operatorname{supp}\hat c\subseteq m_0\Z_n\}$, i.e.\ the
$d_{\max}$-periodic real vectors, of real dimension $d_{\max}$. This count
is read off directly from $|Z_0|=d_{\max}$ and is insensitive to parity:
an alternate derivation via counting conjugate-symmetric pairs
$\{t,-t\bmod n\}$ in $Z_0$ would need to treat the self-paired frequency
$t=n/2$ separately when $n\equiv0\bmod4$ places $n/2\in Z_0$, but that
bookkeeping is specific to the pair-counting route and does not arise
here, since Steps 1--4 characterize $K$ by its Fourier support directly.
\end{proof}

\begin{remark}[Scope of the written derivation]
Steps 2--3 evaluate the same kind of quadratic exponential sum as Step
1's self-transform identity, via the same completion-of-square technique.
The definition of the reality-defect kernel $e$ is the one place this
computation is easy to get wrong: dropping the conjugate term mis-simplifies
the zero set to the wrong condition, structurally the same error as
transferring $\chi$'s self-duality naively (i.e.\ without the affine
correction) to Bj\"orck's construction below (Lemma~\ref{lem:bjaffine}'s
own point of departure). The \emph{conclusion}
$\dim K=d_{\max}$ is independently corroborated: by hand at $n=4$
(self-paired case) and $n=9$ (odd, non-self-paired) via direct enumeration
of $Z_0$, and numerically across all $21+43=64$ tested Zadoff--Chu
configurations (both parities; Section~\ref{sec:zc}'s corollaries below
cite the exact scripts). No counterexample or numerical discrepancy has
been found at any tested configuration.
\end{remark}

\begin{corollary}\label{cor:D1}
Hypothesis (R) holds at every Zadoff--Chu tuple iff $n$ is squarefree ---
in particular at every prime $n$, odd or $n=2$. Theorem~C is therefore
\textbf{unconditional} whenever the Theorem~A limit lies on the symmetry
orbit of a Zadoff--Chu tuple and $n$ is squarefree.
\end{corollary}

\begin{corollary}\label{cor:D2}
At Zadoff--Chu tuples with squarefree $n$ (either parity), the CAZAC
tangent dimension is exactly $(N-1)n+1$: a $1$-dimensional phase circle for
$N=1$ versus an $(n+1)$-dimensional manifold for $N=2$.
\end{corollary}

Numerical verification uses \path{experiments/verify_zc_regularity.py} for
odd lengths
\[
n\in\{3,5,7,9,15,21,25,27,45,49,75\}
\]
($21$ configurations), and
\path{experiments/verify_even_n_regularity.py} for the even lengths
\[
n\in\{2,4,6,8,10,12,14,16,18,20,24,36,50\}
\]
($43$ configurations). All match $\dim K=d_{\max}$ exactly. No
Chebotar\"ev-type nonvanishing-minors
argument is needed: the DFT of a chirp is a chirp, collapsing the
dependency analysis to counting solutions of $t^2\equiv0\bmod n$. Both
scripts are in the companion repository \cite{CazacAlgorithmRepo}.

\begin{remark}[Prior art bearing on the lower bound]
Three older, independent constructions already anticipate part of
Theorem~D --- its \emph{existence} half, that $\dim K$ can be as large as
$d_{\max}$, not just $1$ --- via routes substantially different from, and
more elementary than, this paper's own. \textbf{Backelin (1989)}
\cite{Backelin89}: for $n=m^2\ell$ with $\ell$ squarefree, an explicit
algebraic construction gives an $(m-1)$-parameter family of solutions to
the cyclic-$n$-roots system (the polynomial system underlying CAZAC/
biunimodular sequences), which at $m=d_{\max}$, restricted to unimodular
parameters, matches $d_{\max}-1$ free real phases --- the same count as
$\dim K$ once the always-present global-phase direction is added back.
\textbf{Popovi\'c (1992)} \cite{Popovic92}: modulating a Zadoff--Chu
sequence of length $n=sm^2$ by an \emph{arbitrary} unimodular sequence of
period $m$ again satisfies the ideal periodic autocorrelation property
(their Theorem~1); at $m=d_{\max}$ this gives an explicit
$d_{\max}$-real-parameter family of CAZAC points \emph{through every
Zadoff--Chu tuple} --- confirmed numerically that this family exists
exactly as claimed and has the predicted dimension at $n=4$ ($d_{\max}=2$:
an arbitrary period-$2$ phase pair gives a verified $2$-parameter family
of exact CAZAC points; script: {\tt experiments/verify\_popovic\_n4.py} in
the companion repository \cite{CazacAlgorithmRepo}). \textbf{Bj\"orck--Saffari (1995)} \cite{BjorckSaffari95}:
a further, more general construction (any constant-amplitude length-$N$
seed sequence and any permutation, not just a phase modulation of a fixed
Zadoff--Chu sequence) gives explicit continuous families of CAZAC
sequences of length $N^2$, extended by a further construction (with an
initial parity condition removed by a Kronecker-product trick) to every
non-squarefree length unconditionally. Of the three,
Popovi\'c's is the most directly comparable to this paper's setup (a
family literally through a Zadoff--Chu tuple), and we have \textbf{not}
worked out the precise identification between any of these three explicit
families and this paper's dependency space $K$ (a natural next step; the
dimension matches above are evidence for, not proof of, such an
identification). What Theorem~D adds beyond all three: the \emph{exact}
count (no more than $d_{\max}$, not merely at least $d_{\max}$), at every
Zadoff--Chu tuple and both parities; the reusable chirp-DFT-closure
mechanism itself, which extends to Bj\"orck's non-Zadoff--Chu construction
and to ruling out cubic-or-higher phase; and the connection to hypothesis
(R) and Theorem~C's convergence rate, none of which the 1989--1995
constructions address. Separately: Bj\"orck and Saffari also
\emph{conjectured} the complementary finiteness statement --- that at
squarefree $n$ the full CAZAC variety (not just the Zadoff--Chu orbit) is
finite up to phase --- which remains open in general (settled only for
prime $n$, by Haagerup \cite{Haagerup08}); this is the natural generalization of exactly
the gap Q4 below leaves open (regularity at non-Zadoff--Chu CAZAC points),
not a new question invented by this paper.
\end{remark}

\subsection{The chirp-argument schema, and where it does and does not generalize}

Stripped of Zadoff--Chu-specific notation, Theorem~D's proof needs exactly:
(1) a closure property, $\widehat{x_u}=G_u\cdot x_u\circ\phi_u$ for a
scalar $G_u$ and an index-set automorphism $\phi_u$; (2) an algebraic
addition law for the phase making the key substitution exact; (3) applying
the closure twice, once to reduce the dependency condition to a
correlation equation, once more to diagonalize that equation into a
congruence on the Fourier support of the kernel. This replaces a
linear-algebra rank question with a congruence-counting question, and is a
genuinely reusable schema, not a one-off computation.

\begin{proposition}[Positive lead --- Proved classically, newly connected here]
The Legendre symbol $\chi$ modulo an odd prime $p$ satisfies
$\hat\chi(k)=\chi(k)\,g(\chi)$ for $k\ne0$ ($g(\chi)$ the Gauss sum,
$|g(\chi)|^2=p$): $\chi$ is, up to a scalar, its own DFT, the same closure
shape used by Theorem~D, with trivial reindexing.
\end{proposition}

\begin{proposition}[Obstruction to this specific route --- not a general
impossibility result]
For cubic and higher-degree phase (``algebraic chirp'') sequences,
Theorem~D's specific closure mechanism --- completing the phase by a
shift, the way Step~1 completes the square --- fails, for the same
reason completing the square is special to degree $2$: the addition law
$(m+s)^d=m^d+dsm^{d-1}+\binom d2s^2m^{d-2}+\dots+s^d$ has a mixing term
$dsm^{d-1}$ of degree $d-1$ in $m$, which is linear in $m$ (absorbable into
a pure reindexing of the phase, exactly as Step~1 does) only when $d=2$;
for $d\ge3$ and $\gcd(d,n)=1$, matching the cross term $-km$ (linear in
$m$) forces $ds\equiv0\pmod n$, hence $s\equiv0$, so no nontrivial shift
completes the $d$-th power the way Step~1 completes the square.
Separately, the specific quadratic Gauss sum $G_u$ entering Theorem~D's
Step~1 has magnitude exactly $\sqrt n$, forced by Parseval applied to the
Zadoff--Chu sequence itself (not by a general closed-form evaluation:
classical quadratic Gauss sums have magnitude $\sqrt n$ for $n$ odd, $0$
for $n\equiv2\bmod4$, and $\sqrt{2n}$ for $n\equiv0\bmod4$); Weyl sums of
degree $\ge3$ have no comparable exact evaluation, only upper bounds, so
even if some other closure map existed it would not obviously carry the
exact-modulus property Step~1 needs.
\textbf{Together these show the specific self-transform mechanism used in
Theorem~D's proof does not extend to cubic or higher phase}; this is an
obstruction to \emph{this} route, not a proof that no route exists, and
the mixing-term argument itself is stated only for $\gcd(d,n)=1$ ---
when $d$ and $n$ share a factor, $ds\equiv0\pmod n$ can have nontrivial
solutions and the argument as given does not cover that case.
\end{proposition}

\subsection{The Bj\"orck case: a partial result, not a proof}

Björck's construction (odd prime $p$, Legendre symbol $\chi$)
\cite{Bjorck90} defines $b(m)=e^{i\theta(m)}$ with $\theta$ a two-valued
function of $\chi(m)$, differing by $p\bmod4$. Unlike $\chi$ itself, $b$
is not literally self-dual under the DFT, but a different closure fact
rescues part of the argument.

\begin{lemma}[Proved, exact symbolic computation]\label{lem:bjaffine}
$b$ is exactly \textbf{affine} in $\chi$: $b=A\cdot\mathbf1+B\cdot\chi+C\cdot\delta_0$
with $C=1-A$, where
\[
\begin{aligned}
p\equiv1\pmod4:\quad
 &A=\cos\alpha,\quad B=i\sin\alpha,\\
 &\alpha=\arccos\frac1{1+\sqrt p};\\
p\equiv3\pmod4:\quad
 &A=\frac{1+e^{i\beta}}2,\quad B=\frac{1-e^{i\beta}}2=C,\\
 &\beta=\arccos\frac{1-p}{1+p}.
\end{aligned}
\]
The $3$-dimensional space $T:=\operatorname{span}_\C\{\mathbf1,\chi,\delta_0\}$
closes under the DFT: $\hat b=C\cdot\mathbf1+Bg\cdot\chi+Ap\cdot\delta_0$
(using $F(\mathbf1)=p\delta_0$, $F(\chi)=g\chi$, $F(\delta_0)=\mathbf1$,
$g=g(\chi)$ the Gauss sum, $g=\sqrt p$ for $p\equiv1\pmod4$ and $g=-i\sqrt p$
for $p\equiv3\pmod4$, with the sign fixed by this paper's DFT convention
$\hat x(k)=\sum_mx(m)\omega^{-km}$, i.e.\ $g=\overline{G(\chi)}$ for the
classical Gauss sum $G(\chi)=\sum_m\chi(m)\omega^m=i\sqrt p$ at
$p\equiv3\pmod4$).
\end{lemma}

\begin{lemma}[Proved]\label{lem:bjequiv}
Let $G$ be the group of quadratic-residue units mod $p$ (order $q=(p-1)/2$),
acting by decimation. $b,\hat b$ are pointwise $G$-invariant, so $K$ is a
$G$-invariant subspace of $\R^p$. $\R^p\cong\mathbf1_{\{0\}}\oplus\R[G]\oplus\R[G]$
as a $G$-module, with trivial-isotypic piece exactly $T\cap\R^p$
(dimension $3$).
\end{lemma}

\begin{theorem}[Proved, both residue classes]\label{thm:E0}
$K\cap T=\operatorname{span}\{\mathbf1\}$; i.e.\ restricted to the
$3$-dimensional symmetry-invariant candidate space, the only dependency is
the Parseval one, for every prime $p$.
\end{theorem}
\begin{proof}
Restricting the dependency condition
$\overline{b(m)}\,(F^*(c\odot\hat b))(m)\in\R$ for all $m$ to
$c=\alpha\mathbf1+\beta\chi+\gamma\delta_0\in T$ collapses --- since
$b,\hat b\in T$ closes under pointwise product and $F$ --- to exactly
three scalar equations, one evaluation each at $m=0$, $m\in$~QR,
$m\in$~QNR (the condition depends on $m$ only through $\chi(m)\in\{0,1,-1\}$);
$\alpha$ never appears in any of the three (consistent with
$\mathbf1\in K$ unconditionally), so only the $\beta,\gamma$ coefficients
are at stake. Use $\alpha_0:=\alpha$ and $\beta_0:=\beta$, the same angles
fixed in Lemma~\ref{lem:bjaffine} (not a phase of $Bg$: for
$p\equiv1\pmod4$, $Bg=i\sin\alpha\cdot\sqrt p$ is purely imaginary, never
real, so no polar-decomposition phase of $Bg$ itself is what is meant
here).

For $p\equiv1\pmod4$, eliminating the $m=0$ equation (which involves only
$\beta$) leaves:
$$s\sqrt p\,\beta=0,\qquad -s\gamma\bigl(1+c(p-1)\bigr)=0\quad(c=\cos\alpha_0=\tfrac1{1+\sqrt p},\,s=\sin\alpha_0).$$
$\alpha_0=\arccos\frac1{1+\sqrt p}\in(0,\pi/2)$ since
$\frac1{1+\sqrt p}\in(0,1)$, so $s=\sin\alpha_0>0$; $s\sqrt p\ne0$ forces
$\beta=0$ directly. Also $c=\cos\alpha_0=\frac1{1+\sqrt p}\in(0,1)$
exactly (not merely bounded, and not $\pm1$), so $1+c(p-1)>1>0$ for every
$p\ge5$ (both terms of $c(p-1)$ positive), hence $1+c(p-1)\ne0$ forces
$\gamma=0$.

For $p\equiv3\pmod4$, the QR and QNR equations (the $m=0$ equation is
linearly dependent on these two) are:
$$\sqrt p(1-\cos\beta_0)\,\beta+\tfrac{p-1}2\sin\beta_0\,\gamma=0,\qquad
\tfrac{p+1}2\sin\beta_0\,\gamma=0.$$
$\beta_0=\arccos\frac{1-p}{1+p}$ with $\frac{1-p}{1+p}\in(-1,1)$ strictly
for every $p\ge3$, so $\beta_0\in(0,\pi)$ and $\sin\beta_0>0$; the second
equation then forces $\gamma=0$ directly, since $\tfrac{p+1}2\sin\beta_0\ne0$;
substituting into the first, $\sqrt p(1-\cos\beta_0)\beta=0$ forces
$\beta=0$ as well, since $1-\cos\beta_0\ne0$ ($\beta_0\ne0$).

In both residue classes the system forces $\beta=\gamma=0$ with $\alpha$
free, giving $K\cap T=\operatorname{span}\{\mathbf1\}$ exactly.
\end{proof}

\begin{remark}[What is proved and what remains open]
Lemma~\ref{lem:bjaffine} and Lemma~\ref{lem:bjequiv} and
Theorem~\ref{thm:E0} are complete proofs. What remains \textbf{unproved} is
whether $K$ meets any of the $q-1$ nontrivial isotypic pieces of $\R^p$
under $G$: by Schur's lemma the dependency operator restricted to each
such (complex $2$-dimensional) block is a $2\times2$ matrix, structurally a
Jacobi-sum-type bilinear expression, and its nonvanishing was not proved.
For $p\equiv1\pmod4$, an exact closed form for the block (in terms of
twisted Gaussian periods $\rho_R(k),\rho_N(k)$ and sequence-specific
factors $u,v,w,p_f$) gives $\det=-\tfrac14[e(-k/q)\rho_R(k)^2u(k)v(k)-\rho_N(k)^2p_f(k)w(k)]$;
$\rho_R(k)=\tfrac12[g(\lambda_k)+g(\lambda_k\chi)]$,
$\rho_N(k)=e(-k/2q)\tfrac12[g(\lambda_k)-g(\lambda_k\chi)]$ for
$\lambda_k$ the order-$(p-1)$ character with $\lambda_k(g_0)=e(k/(p-1))$
($g_0$ a primitive root mod $p$). \textbf{One route to $\det=0$ is now
ruled out unconditionally}: the route where $\rho_R(k)$ \emph{and}
$\rho_N(k)$ vanish simultaneously (which would zero out the entire
$2\times2$ block, hence $\det$). $\rho_R(k)=0$ or $\rho_N(k)=0$
individually requires $g(\lambda_k)=\pm g(\lambda_k\chi)$, and the
classical Gauss--Jacobi identity $g(\psi)g(\chi)=J(\psi,\chi)g(\psi\chi)$
re-expresses this as $J(\lambda_k,\chi)=\pm g(\chi)$; since
$J(\lambda_k,\chi)\in\Q(\zeta_{p-1})$ (a $\Z$-linear combination of
$(p-1)$-th roots of unity) while $g(\chi)\in\Q(\zeta_p)\setminus\Q$
(classically $\pm\sqrt p$ or $\pm i\sqrt p$, never rational) and
$\Q(\zeta_{p-1})\cap\Q(\zeta_p)=\Q$ (since $\gcd(p-1,p)=1$), the equality
$J(\lambda_k,\chi)=\pm g(\chi)$ would force $g(\chi)\in\Q$, a
contradiction. So $\rho_R(k),\rho_N(k)\ne0$ for every prime and every
nontrivial $k$, \textbf{proved}, not merely checked (script: {\tt
experiments/verify\_bjorck\_jacobi\_unification.py} in the companion
repository \cite{CazacAlgorithmRepo}) --- a clean standalone fact, but
\textbf{one that does not, by itself, narrow the determinant question}:
neither $\rho_R(k)=0$ nor $\rho_N(k)=0$ was ever a necessary or a
sufficient condition for $\det=0$ on its own (e.g.\ at $\rho_R(k)=0$,
$\det$ collapses to $\tfrac14\rho_N(k)^2p_f(k)w(k)$, itself zero only if
$p_f$ or $w$ also vanishes), so this proof does not directly rule out
$\det=0$ in general. $\det$ is a \emph{difference} of two terms, and the
actual open condition is the cross-term coincidence
$e(-k/q)\rho_R(k)^2u(k)v(k)=\rho_N(k)^2p_f(k)w(k)$, which can hold even
with $\rho_R,\rho_N,u,v,w,p_f$ all individually nonzero.
$u(k)v(k)-w(k)p_f(k)$ has an explicit nonzero closed form
($-16\,A\,\varepsilon\,g\,C\sin^2\alpha$ for the constants $A,C$ of
Lemma~\ref{lem:bjaffine} and $\varepsilon:=(-1)^k$; itself never zero,
since $A=\cos\alpha\in(0,1)$, $C=1-A\ne0$, $g=\sqrt p\ne0$, and
$\sin\alpha\ne0$), confirming the
cross-term is a genuinely separate condition, not reducible to $\rho_R,
\rho_N$ vanishing. Checked directly, not assumed away, \textbf{for
$p\equiv1\pmod4$ in $[5,60)$}: at the $k$ minimizing $|\det|/p^2$ for each
such prime, $\min(|\rho_R(k)|,|\rho_N(k)|)/\sqrt p$ stays in $[0.11,0.6]$
--- \textbf{not small} --- even as $|\det|/p^2$ itself shrinks by three
orders of magnitude, confirming that in \emph{this} range the
near-degeneracy traces to the cross-term, not to $\rho_R$ or $\rho_N$
(script: {\tt experiments/verify\_bjorck\_crossterm\_gap.py}, same
repository). \textbf{This does not extend to $p=71$}: that prime is
$\equiv3\pmod4$, outside both the closed form above and this sweep's
range, and at $p=71,k=17$ specifically $\min(|\rho_R|,|\rho_N|)/\sqrt p
\approx5.3\times10^{-7}$ --- there, unlike the $p\equiv1\pmod4$ range
just checked, the near-degeneracy traces directly to $\rho_R$ or $\rho_N$
being small, not to the cross-term; the two residue classes are not
shown to behave the same way, and the $p\equiv3\pmod4$ closed form
needed to even pose the cross-term question there has not been derived.
The cross-term coincidence remains open and unresolved for
$p\equiv1\pmod4$. Separately, the
relevant determinants are confirmed numerically nonzero at every one of
$44$ tested primes $p\in[5,200)$ (both residue classes, $\dim K=1$
throughout) --- at $p=71$ the gap is a factor of $\approx4\times10^3$
smaller than at the next-nearest tested prime, a genuine near-degeneracy,
independently confirmed not to be a floating-point artifact by a
$60$-digit-precision (mpmath) recomputation of this specific block (top
singular value matches to $\sim10^{-15}$ relative precision, the
near-zero singular value moves by only $\sim1\%$ of its own scale between
float64 and $60$-digit precision, not by orders of magnitude). Script:
{\tt experiments/verify\_bjorck\_isotypic.py} in the companion repository
\cite{CazacAlgorithmRepo} (the $44$-prime determinant check; a
separately-run set of $14$ primes, all contained in this same tested
range, is used for the affine-closure check of Lemma~\ref{lem:bjaffine},
in {\tt experiments/verify\_bjorck\_regularity.py}).
\textbf{The overall regularity claim at Bj\"orck points is not proved in
this paper}: what is established here is a full proof restricted to the
symmetry-invariant subspace, an unconditional proof ruling out one route
(both twisted periods vanishing at once) by which the remaining
determinant could fail, and strong (but non-uniform) numerical evidence,
not a proof, for the complementary directions as a whole. Full regularity
--- for every prime and both residue classes --- is proved in the
literature by an independent argument of the same general character
(decimation-group isotypic decomposition, a Gauss-sum ratio shown never
to be a root of unity via Stickelberger's theorem): Benoist
\cite{Benoist24}, Prop.~A.3. This paper's own contribution regarding
Bj\"orck's construction is the connection to hypothesis (R) and
Theorem~C's convergence rate, not the underlying regularity fact itself.
\end{remark}

\section{Open problems}\label{sec:open}

\noindent\textbf{Q1 (settled).} Does $x^\ell$ converge to a single point? Yes ---
Theorem~\ref{thm:A}.

\noindent\textbf{Q2 (open, reframed).} Characterize all fixed points of $(\star)$
for general $n,N$. The stationarity system is a KKT eigenvalue equation
$x^*=(\Lambda+F^*MF)x^*$ for diagonal multipliers, structurally reminiscent
of discrete prolate/Slepian time-and-band-limiting operators. Complete for
$n=2$ (four symmetry classes); open in general. A sharper sub-question:
are non-CAZAC fixed points always measure-zero repellers, or can they
attract an open set of initializations?

\noindent\textbf{Q3 (partially answered).} Why does $N=2$ empirically always
succeed where $N=1$ sometimes fails? Corollary~\ref{cor:D2}'s exact
dimension count, $\dim\CAZ(n,N)=(N-1)n+1$ at squarefree Zadoff--Chu
tuples, suggests a mechanism (more coupled sequences means more escape
directions), though this is proved only at those specific points, not as
a general lower bound at every CAZAC point (which would additionally
require $\CAZ(n,N)$ to be a manifold there, not established in general);
the pathological fixed-point classes are the same for $N=1,2$, so the gap
is about basin measure, not existence, and remains open.

\noindent\textbf{Q4 (essentially settled locally; one remaining gap).} Rate of
convergence near $\CAZ(n,N)$: settled by Theorem~C at regular points, and
Theorem~D now makes regularity unconditional at Zadoff--Chu points of
either parity for squarefree $n$. Remaining: (R) at non-Zadoff--Chu CAZAC
points in general is open in this paper (Bj\"orck's own construction is
resolved in the literature, \cite{Benoist24}, though not by this paper's
methods); $n$-uniformity of
the true sharp-growth constant is conjectural (the proved bound scales
like $c_N/n$; whether the true constant does too is open and not backed
here by a specific numerical check). This remaining gap is not a new
question: it is the natural generalization of Bj\"orck and Saffari's own
still-open finiteness conjecture (\S\ref{sec:zc}, remark after
Corollary~\ref{cor:D2}) --- settled only at prime $n$, by Haagerup \cite{Haagerup08}.

\noindent\textbf{Q5 (open, deprioritized).} What, if anything, do the
preliminary solves buy? The reference implementation solves successively at
$\tau=1,1/2,\ldots,1/10$ before resetting to $\tau=1$. Their effect has not
been tested by ablation, and no theoretical benefit is claimed.

\noindent\textbf{Q6 (the central open question).} For $N\ge2$ and Gaussian
initialization, does the iteration converge to $\CAZ(n,N)$ with
probability $1$ (or $\ge1-c^n$)? Requires a full fixed-point census (Q2)
and a basin/measure argument. No tracked experiment in this repository
currently estimates that success probability for $(\star)$. The
$15/50$ success count at $n=167$ printed by the tracked projection script
belongs to a different IPUC-style baseline and is not evidence for this
question.

\noindent\textbf{Q7 (partially settled by citation, residual part open).}
Does every trajectory of the elliptope iteration converge to a rank-1
point? Theorem~\ref{thm:ell} proves single-point convergence
unconditionally but not convergence to an extreme point specifically.
The \emph{pointwise} half of this question is already settled ---
Remark~\ref{rem:fkp-thm20} cites \cite{FKP21}'s own Theorem~20: rank-1
points are exactly the fixed points whose full neighborhood converges to
them, and every other fixed point admits an explicit escaping direction.
What remains open is the \emph{measure-theoretic} refinement: does the
set of initializations whose trajectory converges to a non-vertex fixed
point have Gaussian measure zero? \cite{FKP21}'s own escaping-curve
argument (Prop.~19, Remark~\ref{rem:fkp-thm20}) only produces one
escaping point per neighborhood, not universal escape, so it leaves open
whether non-vertex elliptope points are repelling outright or could
instead be collectively (not individually) attracting, the way
\cite{FKP21}'s unrelated cone example (\S3, Example~5) shows can happen
in general for this style of iteration. This residual
question is structurally the same conjecture as Q6 in a second, unrelated
domain: if the measure is zero, it is independent evidence for Q6's
general shape; if not, CAZAC's extra structure (the dimension count of
Corollary~\ref{cor:D2}) may be doing more work than currently credited.
The numerical evidence (Prop.~\ref{prop:ell4}, $140$ trials across
$n\in\{3,\dots,20\}$) shows zero exceptions, consistent with --- but
not a proof of --- the measure-zero reading. \emph{Attack plan:} a
basin/measure argument analogous to Q6's, using the pointwise dichotomy
of Theorem~20 as the starting classification of which points can even be
candidates for a positive-measure limit.

\section{Verification status and provenance}

\textbf{Tiers.} \emph{Unconditional}: complete proof, no open hypothesis.
\emph{Proved conditional on (R)}: complete proof, resting on regularity
hypothesis (R). \emph{Proved conditional on (R) and an open sub-case}:
complete proof modulo one explicitly identified, unresolved lemma.
\emph{Verified numerically}: no proof; a script reproduces the claim to a
stated tolerance. \emph{Conjecture}: stated as a target, not established
either way.

\textbf{Principal claims.}
\begin{itemize}
\item Propositions~\ref{prop:1}--\ref{prop:4} (generalized power method,
max-norm = CAZAC, monotone ascent, limit points fixed): \emph{Unconditional}.
\item Theorem~A (global single-point convergence): \emph{Unconditional}
as a mathematical claim, but its descent-argument mechanism is a special
case of Akta\c{s}--Kroer's \cite{AK25} general Frank-Wolfe/KL framework
(2025); the paper's own contribution here is the CAZAC-to-$\max\|x\|^2$
reduction (Prop.~\ref{prop:2}) that makes that template apply, not the
descent argument itself --- see the note at Theorem~A.
\item Elliptope semialgebraicity and H1--H3 transfer
(Props.~\ref{prop:ell1}--\ref{prop:ell2}): \emph{Unconditional}, every $n$.
Fixed-point continuum (Prop.~\ref{prop:ell3}): \emph{Unconditional, $n>3$
only} (infinitude by citation to \cite{FKP21}, positive-dimensionality by
semialgebraicity, plus a machine-exact illustrative identity at $I_n$); at
$n\le3$ the fixed-point set is finite (per \cite{FKP21}'s own count), so
the statement is restricted to $n>3$, not vacuously true for all $n$.
Single-point convergence (Theorem~\ref{thm:ell}): \emph{Unconditional for
every $n$} --- this is a proved theorem, not a numerical finding; the
``despite an infinite continuum'' framing is specific to $n>3$, but
convergence itself does not depend on that framing. The pointwise
vertex/non-vertex attractive
dichotomy (Remark~\ref{rem:fkp-thm20}): \emph{Unconditional}, by direct
citation to \cite{FKP21} Theorem~20 (not a new argument here).
\emph{Separately}, that generic trajectories escape onto a rank-1 point
specifically (Prop.~\ref{prop:ell4}): \emph{Verified numerically only}
($140$ trials, $n\in\{3,\dots,20\}$) --- this is not proved. The residual,
measure-zero form of universal rank-1 convergence (Q7): \emph{Conjecture}.
\item Theorem B (obstruction, exposed non-CAZAC fixed point):
\emph{Unconditional} for every $N$ (part (a)'s statement and proof cover
general $N$: the coupled spectral condition $S_x(k)=Nn$ on $K$ is pinned,
individual per-sequence moduli are not, and only $N=1$ collapses to a
per-block modulus condition).
\item Theorem C, parts (a),(b),(c), including the full contradiction-based
proof of part (c): \emph{Proved conditional on hypothesis (R)}; all three
parts carry the same status tag.
\item Theorem D (regularity of Zadoff--Chu points, $\dim K=d_{\max}$),
\textbf{both odd and even $n$}: \emph{Unconditional}, a complete proof for
all $n$ regardless of parity, with Steps 2--3 spelled out via the
reality-defect kernel $e$ (see the ``Scope of the written derivation''
remark after the proof) --- independently corroborated by hand at $n=4,9$
and numerically at all 64 tested configurations. Prior art disclosed
(remark after Corollary~\ref{cor:D2}): three older, independent
constructions --- Backelin (1989) \cite{Backelin89}, Popovi\'c (1992)
\cite{Popovic92}, Bj\"orck--Saffari (1995) \cite{BjorckSaffari95} --- each
give, by unrelated elementary methods, explicit families matching the
\emph{existence} half of this count ($\dim K$ can reach $d_{\max}$) at
non-squarefree $n$; Theorem~D's own contribution is the exact count and
the reusable mechanism, not the existence of degeneracy per se.
Corollaries~\ref{cor:D1}--\ref{cor:D2}: \emph{Unconditional} (Corollary~1
conditionally activates Theorem~C, but is itself unconditional as a
dimension/regularity statement).
\item Legendre-symbol closure positive lead: classical fact, newly
connected to the schema, \emph{Unconditional} as stated (no claim of
Bj\"orck regularity is made from it alone). Cubic/higher-phase obstruction:
\emph{proved for the specific completing-the-square mechanism} used in
Theorem~D's proof, when $\gcd(d,n)=1$ (the mixing term becomes nonlinear
in $m$ for $d\ge3$); \textbf{not} a general impossibility result, and not
addressed when $\gcd(d,n)>1$.
\item Bj\"orck partial result: Lemma~\ref{lem:bjaffine},
Lemma~\ref{lem:bjequiv}, Theorem~\ref{thm:E0} ($K\cap T=\operatorname{span}\{\mathbf1\}$):
\emph{Unconditional}. Within the residual (non-symmetry-invariant)
directions: $\rho_R(k),\rho_N(k)\ne0$ for every prime and every nontrivial
$k$ is \emph{Unconditional} (proved, an elementary cyclotomic-field
argument, not merely checked --- \S\ref{sec:zc}, Björck subsection), but
neither $\rho_R(k)=0$ nor $\rho_N(k)=0$ was ever a necessary or a
sufficient condition for the full isotypic-block determinant to vanish
on its own --- this rules out only the sub-case where both vanish at
once, not the actual cross-term condition, which this paper's own methods
do not resolve (verified numerically only, at 44 primes, $p\in[5,200)$).
Full regularity (R) at Bj\"orck points is \textbf{not proved in this
paper}; it is available in the literature \cite{Benoist24}, via a similar
Gauss-sum technique, for every prime and both residue classes.
\item Q1--Q7: precise open-problem statements as in Section~\ref{sec:open}; Q1 is
settled (restated as a theorem above), the rest range from ``partially
answered'' to fully open conjectures, as stated there.
\end{itemize}

\textbf{Provenance.} This document was written with substantial assistance
from large language models (Claude, Anthropic) under the direction of
Jeffery Kline, drawing on a longer working record of derivations kept in
the companion repository \cite{CazacAlgorithmRepo}'s {\tt theory/}
directory. Numerical claims name their supporting scripts in that
repository's {\tt experiments/} directory. Claims whose supporting companion
artifact is unavailable are identified as such in the repository's release
documentation.

\begin{thebibliography}{99}

\bibitem{ABS13} H. Attouch, J. Bolte, B. F. Svaiter, \emph{Convergence of
descent methods for semi-algebraic and tame problems: proximal algorithms,
forward--backward splitting, and regularized Gauss--Seidel methods},
Math. Program. \textbf{137} (2013), 91--129.

\bibitem{BDLS07} J. Bolte, A. Daniilidis, A. S. Lewis, M. Shiota,
\emph{Clarke subgradients of stratifiable functions}, SIAM J. Optim.
\textbf{18}(2) (2007), 556--572. arXiv:math/0601530.

\bibitem{FKP21} P. Felzenszwalb, C. Klivans, A. Paul, \emph{Iterated Linear
Optimization}, Quart. Appl. Math. \textbf{79}(4) (2021), 601--615.
arXiv:2012.02213.

\bibitem{AK25} F. S. Akta\c{s}, C. Kroer, \emph{Strongly Convex Maximization
via the Frank-Wolfe Algorithm with the Kurdyka-{\L}ojasiewicz Inequality},
arXiv:2505.00221 (2025).

\bibitem{JNRS10} M. Journ\'ee, Y. Nesterov, P. Richt\'arik, R. Sepulchre,
\emph{Generalized Power Method for Sparse Principal Component Analysis}, J. Mach. Learn. Res.
\textbf{11} (2010), 517--553. arXiv:0811.4724.

\bibitem{LLM09} A. S. Lewis, D. R. Luke, J. Malick, \emph{Local linear
convergence for alternating and averaged nonconvex projections}, Found.
Comput. Math. \textbf{9}(4) (2009), 485--513. arXiv:0709.0109.

\bibitem{DL12} D. Drusvyatskiy, A. S. Lewis, \emph{Optimality,
identifiability, and sensitivity}, Math. Program. \textbf{147} (2014),
467--498. arXiv:1207.6628.

\bibitem{Chu72} D. C. Chu, \emph{Polyphase codes with good periodic
correlation properties}, IEEE Trans. Inform. Theory \textbf{18}(4) (1972),
531--532.

\bibitem{Amis25} K. Amis, E. Boutillon, E. Boutillon, \emph{CAZAC sequence
generation of any length with iterative projection onto unit circle:
principle and first results}, arXiv:2509.05097 (2025).

\bibitem{Backelin89} J. Backelin, \emph{Square multiples $n$ give
infinitely many cyclic $n$-roots}, Reports, Matematiska Institutionen
\textbf{8}, Stockholms Universitet, 1989.

\bibitem{Popovic92} B. M. Popovi\'c, \emph{Generalized chirp-like
polyphase sequences with optimum correlation properties}, IEEE Trans.
Inform. Theory \textbf{38}(4) (1992), 1406--1409.

\bibitem{Bjorck90} G. Bj\"orck, \emph{Functions of modulus $1$ on $\Z_n$
whose Fourier transforms have constant modulus, and ``cyclic $n$-roots''},
in: Recent Advances in Fourier Analysis and its Applications, NATO ASI
Series C \textbf{315}, Kluwer, 1990, 131--140.

\bibitem{BjorckSaffari95} G. Bj\"orck, B. Saffari, \emph{New classes of
finite unimodular sequences with unimodular Fourier transforms. Circulant
Hadamard matrices with complex entries}, C. R. Acad. Sci. Paris S\'er. I
\textbf{320}(3) (1995), 319--324.

\bibitem{Benoist24} Y. Benoist, \emph{Fourier transform in cyclic groups},
arXiv:2406.11529 (2024).

\bibitem{Haagerup08} U. Haagerup, \emph{Cyclic $p$-roots of prime lengths
$p$ and related complex Hadamard matrices}, arXiv:0803.2629 (2008).

\bibitem{CazacAlgorithmRepo} J. Kline, \emph{cazac-algorithm}: companion
research repository holding the {\tt theory/} derivations, the dated
{\tt AUDIT-LEDGER.md}, and the {\tt experiments/} scripts (including
{\tt verify\_zc\_regularity.py}, {\tt verify\_even\_n\_regularity.py},
{\tt verify\_elliptope\_convergence.py}, {\tt verify\_elliptope\_rank\_sweep.py},
{\tt verify\_bjorck\_regularity.py}, {\tt verify\_bjorck\_isotypic.py},
{\tt verify\_bjorck\_jacobi\_unification.py},
{\tt verify\_bjorck\_crossterm\_gap.py}, {\tt verify\_popovic\_n4.py}), and
a {\tt docs/} directory holding an earlier paper draft that backs one
figure (Q6) by direct citation rather than a script; together these
reproduce every numerical claim in this paper, 2026. \textbf{Not publicly
hosted at time of writing} --- no location is asserted here because none
can currently be given; a reader cannot yet independently fetch this
repository.

\end{thebibliography}

\end{document}
